\section{4. Model Formulation}
To evaluate the stability regimes induced by status-based peer sanctioning, this section formalizes a discrete-time multi-agent consensus model\cite{degroot1974reaching}. The framework extends the classical DeGroot averaging protocol by coupling two distinct behavioral mechanisms: an endogenous, state-dependent suppression of outgoing influence targeting positive outliers, and an exogenous cognitive anchor representing partial stubbornness\cite{degroot1974reaching}. The mathematical formulation defines the baseline interaction topology, the relative status calculation, the functional forms of the suppression penalty, dynamic row normalization, cognitive anchoring, and localized kinship protection\cite{degroot1974reaching}.

\subsection{Network Initialization and the Baseline Influence Matrix}
Let $G = (V, E)$ represent a directed social network comprising $N$ autonomous agents, with node set $V = {1, 2, \dots, N}$ and edge set $E \subseteq V \times V$\cite{degroot1974reaching}. The structural topology is encoded by an adjacency matrix $A = [a_{ij}] \in \{0,1\}^{N \times N}$, where $a_{ij} = 1$ indicates that agent $i$ can receive information directly from agent $j$, and $a_{ij} = 0$ otherwise\cite{degroot1974reaching}.

To establish an unpenalized, egalitarian reference state and prevent structural bipartite oscillations, every agent is assigned a baseline self-loop: $$a_{ii} = 1, \quad \forall i \in V$$ Each agent $i$ is initialized with a private scalar estimate $x_i(0) \in \mathbb{R}$\cite{degroot1974reaching}. The baseline influence matrix $W(0) = [w_{ij}(0)] \in \mathbb{R}^{N \times N}$ is constructed by normalizing the adjacency matrix $A$ across rows\cite{degroot1974reaching}: $$w_{ij}(0) = \frac{a_{ij}}{\sum_{k=1}^N a_{ik}}$$ Because each row of $A$ contains at least its self-loop ($\sum_k a_{ik} \ge 1$), the baseline matrix $W(0)$ is strictly row-stochastic, representing an egalitarian interaction network where each agent distributes equal trust across all active communication channels\cite{degroot1974reaching}.

\subsection{State-Dependent Deviation and Status Calculation}
In sociological leveling dynamics, peer sanctioning is not governed by symmetric opinion distance (homophily), but by an individual's vertical status relative to the collective distribution\cite{degroot1974reaching}. This relative status is evaluated at each discrete time step $t$ by calculating the standardized deviation of each agent's estimate from the unweighted population mean\cite{degroot1974reaching}.

Let $\mu(t)$ denote the unweighted network mean and $\sigma(t)$ denote the population standard deviation of the state distribution at time $t$\cite{degroot1974reaching}: $$\mu(t) = \frac{1}{N} \sum_{i=1}^N x_i(t)$$ $$\sigma(t) = \sqrt{\frac{1}{N} \sum_{i=1}^N \big(x_i(t) - \mu(t)\big)^2}$$

The standardized status deviation $D_i(t)$ for agent $i$ is defined as\cite{degroot1974reaching}: $$D_i(t) = \frac{x_i(t) - \mu(t)}{\sigma(t) + \eta}$$ where $\eta > 0$ is a small positive regularization constant (fixed at $\eta = 10^{-8}$ in numerical implementations)\cite{acemoglu2013opinion}. The parameter $\eta$ resolves the scale-invariant singularity that occurs when the network approaches consensus ($\sigma(t) \to 0$)\cite{acemoglu2013opinion,hajnal1958weak}. Without regularization, infinitesimal numerical perturbations around consensus cause $D_i(t)$ to fluctuate across finite non-zero values, inducing artificial high-frequency chattering in the suppression factor\cite{acemoglu2013opinion,hajnal1958weak}. The constant $\eta$ guarantees uniform continuity across the consensus manifold, ensuring that $D_i(t) \to 0$ as states contract\cite{acemoglu2013opinion,hajnal1958weak}.

The standardized score $D_i(t)$ provides a dimensionless metric of relative outperformance\cite{degroot1974reaching}. If $D_i(t) > 0$, agent $i$'s estimate exceeds the group average, designating the agent as a "tall poppy" subject to peer sanctioning\cite{degroot1974reaching}. If $D_i(t) \le 0$, the agent's estimate resides at or below the collective average\cite{degroot1974reaching}.

\subsection{The Outlier-Suppression Mechanism}
Peer sanctioning is formalized through a state-dependent suppression factor $f_i(t) \in [0, 1]$, which dictates the fraction of agent $i$'s outgoing broadcast influence that survives group leveling\cite{degroot1974reaching}.

The suppression operator is strictly asymmetric and directional\cite{degroot1974reaching}. Agents sitting at or below the collective average experience no social penalty, preserving full communication throughput\cite{degroot1974reaching}: $$f_i(t) = 1.0, \quad \text{if } D_i(t) \le 0$$ For agents exceeding the collective mean ($D_i(t) > 0$), the suppression factor decreases monotonically with positive deviance\cite{degroot1974reaching}: $$f_i(t) = \max\big(0, 1 - P \cdot h(D_i(t))\big), \quad \text{if } D_i(t) > 0$$ Here, $P \in [0, 1]$ represents the global penalty magnitude, governing the societal severity of peer sanctioning\cite{degroot1974reaching}. When $P = 0$, the network operates as an unpenalized consensus system; when $P = 1$, complete communication severance is enforced against extreme outliers\cite{degroot1974reaching}. The function $h: \mathbb{R}_{\ge 0} \to \mathbb{R}_{\ge 0}$ governs the functional response curve of the leveling penalty\cite{degroot1974reaching}.

\subsection{Dynamic Penalty Function Shapes}
To establish whether ergodicity failure and non-convergent cycling represent general properties of asymmetric suppression rather than artifacts of a specific functional continuity, three distinct mathematical formulations for $h(D)$ are defined\cite{degroot1974reaching}:

Linear Penalty ($h(D) = D$): A smooth, proportional suppression model where peer penalization scales linearly with standardized deviation\cite{degroot1974reaching}. Mild outliers experience marginal attenuation, whereas extreme outliers are driven toward zero influence\cite{degroot1974reaching}: $$f_i(t) = \max(0, 1 - P \cdot D_i(t))$$ Under maximum penalty ($P = 1$), influence is completely severed ($f_i(t) = 0$) for any deviation $D_i(t) \ge 1.0$\cite{acemoglu2013opinion}.

Step Penalty ($h(D) = \mathbb{1}[D > \theta]$): A discontinuous, threshold-activated sanctioning model with a social tolerance boundary fixed at $\theta = 0.5$\cite{degroot1974reaching}: $$f_i(t) = \begin{cases} 1 - P, & \text{if } D_i(t) > 0.5 \\ 1.0, & \text{if } D_i(t) \le 0.5 \end{cases}$$ If an agent's deviation crosses the tolerance threshold, its outgoing influence drops discontinuously by $P$\cite{degroot1974reaching}. This tests the network's resilience to sudden, structural communication fractures\cite{degroot1974reaching}.

Tanh Penalty ($h(D) = \tanh(D)$): A smooth, saturating suppression model where the marginal penalty diminishes for extreme deviations\cite{degroot1974reaching}: $$f_i(t) = \max(0, 1 - P \cdot \tanh(D_i(t)))$$ Because $\tanh(D) < 1$ for all finite $D$, the term $1 - P \cdot \tanh(D)$ remains strictly positive for all finite states whenever $P \le 1$, modeling an environment where an outlier's broadcast weight is heavily attenuated but never completely severed\cite{acemoglu2013opinion,hajnal1958weak}.

\subsection{Dynamic Weight Modification and Row Normalization}
Once the suppression factors $f_j(t)$ are calculated for all agents, the network's influence matrix is reweighted\cite{degroot1974reaching}. If agent $j$ is penalized, every listening peer $i$ scales down the weight assigned to $j$\cite{degroot1974reaching}.

The raw modified weight $w'_{ij}(t)$ is obtained by scaling the baseline edge weight by sender $j$'s current suppression factor\cite{degroot1974reaching}: $$w'_{ij}(t) = w_{ij}(0) \cdot f_j(t)$$

Because this operation scales columns corresponding to penalized agents downward, the intermediate matrix $W'(t) = [w'_{ij}(t)]$ loses its row-stochastic property\cite{degroot1974reaching}. To ensure that agent updates remain valid convex combinations, the matrix is row-normalized at every time step\cite{degroot1974reaching}. The final time-varying influence weight $\tilde{w}_{ij}(t)$ is defined as\cite{degroot1974reaching}: $$\tilde{w}_{ij}(t) = \frac{w'_{ij}(t)}{\sum_{k=1}^N w'_{ik}(t)}$$

If a pathological condition occurs where all incoming neighbors of agent $i$ are extreme outliers whose suppression factors evaluate to zero ($\sum_k w'_{ik}(t) = 0$), agent $i$ falls back to complete self-reliance\cite{degroot1974reaching}: $$\tilde{w}_{ii}(t) = 1.0, \quad \text{and} \quad \tilde{w}_{ij}(t) = 0.0 \quad \forall j \ne i$$

This dynamic normalization guarantees that the instantaneous matrix $\tilde{W}(t) = [\tilde{w}_{ij}(t)]$ strictly satisfies row-stochasticity at every time step $t$\cite{degroot1974reaching}: $$\tilde{w}_{ij}(t) \ge 0 \quad \forall i, j \in V, \qquad \sum_{j=1}^N \tilde{w}_{ij}(t) = 1.0 \quad \forall i \in V$$

While row-stochasticity is strictly preserved, column-stochasticity is systematically broken\cite{hajnal1958weak}. Because the columns of upper outliers are discounted prior to normalization, the column sums satisfy $\sum_{i=1}^N \tilde{w}_{ij}(t) < 1$ for agents above the mean and $\sum_{i=1}^N \tilde{w}_{ik}(t) > 1$ for unpenalized agents below the mean\cite{hajnal1958weak}. This asymmetry induces a persistent downward drift in the collective network mean\cite{acemoglu2013opinion,hajnal1958weak}: $$\mu(t+1) \le \mu(t)$$

\subsection{Cognitive Anchoring and Partial Stubbornness}
Applying the dynamic weight matrix $\tilde{W}(t)$ directly to a standard DeGroot update produces complete assimilation\cite{degroot1974reaching}. If an accurate agent $i$ sits above the mean, its outgoing broadcast weights are suppressed ($\tilde{w}_{ji}(t) \to 0$ for all $j \ne i$)\cite{degroot1974reaching}. The group stops listening to agent $i$ and averages strictly among lower-performing peers, causing the group mean to drift downward\cite{degroot1974reaching}.

However, because the penalty is strictly directional, agent $i$ continues to absorb incoming influence from the rest of the network\cite{degroot1974reaching}. Agent $i$ averages its own state with the descending group mean, causing its estimate to decay\cite{degroot1974reaching}. As agent $i$'s estimate approaches the collective mean, its standardized deviation $D_i(t)$ shrinks toward zero, which lifts the penalty\cite{degroot1974reaching}. By the time the penalty is removed, agent $i$'s accurate information has been permanently erased\cite{degroot1974reaching}. The system collapses into a biased static consensus (Regime 1), failing to model the persistent conflict documented in empirical status leveling\cite{degroot1974reaching}.

To capture the cognitive resistance individuals display when holding private evidence against a conformist group, the model introduces a partial stubbornness parameter, $\gamma_{stub} \in [0, 1)$\cite{degroot1974reaching}. At each discrete step, each agent anchors a fraction $\gamma_{stub}$ of its state to its initial private estimate $x_i(0)$, while exposing the remaining $1 - \gamma_{stub}$ to peer averaging\cite{degroot1974reaching}: $$x_i(t+1) = \gamma_{stub} x_i(0) + (1 - \gamma_{stub}) \sum_{j=1}^N \tilde{w}_{ij}(t) x_j(t)$$

In matrix notation, the global state vector evolves according to the non-homogeneous affine dynamical system\cite{acemoglu2013opinion,hajnal1958weak}: $$x(t+1) = (1 - \gamma_{stub}) \tilde{W}(t) x(t) + \gamma_{stub} x(0)$$

Mathematically, $\gamma_{stub}$ acts as a constant restoring force\cite{degroot1974reaching}. While social influence $(1 - \gamma_{stub}) \tilde{W}(t) x(t)$ continuously pulls the penalized outlier toward the collective mean, the anchor $\gamma_{stub} x_i(0)$ prevents the agent's estimate from dropping below a fixed cognitive floor\cite{degroot1974reaching}. This floor preserves the agent's positive standardized deviation, keeping the leveling penalty active and enabling the parameter space required for non-convergent cycling (Regime 3)\cite{degroot1974reaching,acemoglu2013opinion}. When $\gamma_{stub} = 0$, the system collapses unconditionally into biased consensus\cite{degroot1974reaching,acemoglu2013opinion}.

\subsection{Kinship Clustering and Localized In-Group Immunity}
The baseline formulation assumes status-based penalization operates uniformly across the entire graph\cite{degroot1974reaching}. In human social networks, leveling pressures are frequently mitigated by localized in-groups, mutual defense alliances, or elite coalitions\cite{degroot1974reaching,hajnal1958weak}. To formalize structural immunity against status penalties, a kinship cluster $C \subset V$ of size $|C| = K$ is embedded within the network\cite{degroot1974reaching}.

Within this cluster, agents unconditionally exempt each other from status suppression, modeling mutual in-group protection\cite{degroot1974reaching,acemoglu2013opinion}. For any directed communication link $(j, i)$ from sender $j$ to receiver $i$, the intermediate unnormalized weight is governed by the in-group rule\cite{acemoglu2013opinion}: $$w'_{ij}(t) = \begin{cases} w_{ij}(0), & \text{if } i \in C \text{ and } j \in C \\ w_{ij}(0) \cdot f_j(t), & \text{otherwise} \end{cases}$$

This formulation guarantees that\cite{acemoglu2013opinion}:

Mutual In-Group Exemption: When both agents belong to the kinship cluster ($i, j \in C$), sender $j$ is immune to penalization, and receiver $i$ processes $j$'s broadcast at full baseline weight ($w'_{ij}(t) = w_{ij}(0)$) regardless of $j$'s standardized deviance\cite{degroot1974reaching,acemoglu2013opinion}.

Sanctioning of External Outliers: When a kinship member $i \in C$ receives information from an external outlier $j \notin C$, the penalty factor $f_j(t)$ is actively enforced ($w'_{ij}(t) = w_{ij}(0) f_j(t)$)\cite{acemoglu2013opinion}. Kinship members do not act as naive listeners; they protect their own members while continuing to sanction external peers\cite{acemoglu2013opinion}.

Outward Projection: A kinship member $j \in C$ who is an upper outlier continues to experience suppression from external non-kin listeners $i \notin C$ ($w'_{ij}(t) = w_{ij}(0) f_j(t)$), ensuring that immunity remains localized to in-group edges\cite{acemoglu2013opinion}.

The kinship size parameter $K$ is evaluated at discrete fractions of the network population: $$K \in \left\{ N, \frac{N}{2}, \frac{N}{5}, \frac{N}{10} \right\}$$ When $K = N$, the entire network is mutually exempt, deactivating the tall poppy penalty and reverting the system to a classical Friedkin-Johnsen consensus model\cite{degroot1974reaching}. As $K$ shrinks to a minor fraction of the network, localized protection is gradually overwhelmed by global network suppression\cite{degroot1974reaching}. This parameter quantifies the exact critical percolation threshold at which an unpenalized core preserves global ergodicity and immunizes the collective against non-convergent cycling\cite{degroot1974reaching,hajnal1958weak}.

